\subsection{凸锥}

定义 3. --- 仿射空间 E 的子集 C 称为以 \(x_0\) 为顶点的锥，如果 C 在以 \(x_0\) 为中心、比 \(> 0\) 的所有位似变换下不变。

在本目及下一目中，我们将假设已把所考虑的锥的顶点取为 E 中的原点；即假设 E 是向量空间，并且当我们谈到锥时，应理解该锥以 0 为顶点。形如 \(\lambda a\)（其中 \(\lambda > 0\)，相应地 \(\lambda \geqslant 0\)）的点组成的集（其中 \(a\) 是非零向量）称为以 0 为起点的开半直线（相应地闭半直线）。

以 0 为顶点的锥 C 若满足 \(0 \in C\) 则称为尖锥，否则称为无尖锥。尖锥或者是单点集 \(\{0\}\)，或者是以 0 为起点的一族闭半直线之并。无尖锥是以 0 为起点的开半直线之并（可能为空）。若 C 是无尖锥，则 \(C \cup \{0\}\) 是尖锥。若 C 是尖锥，则 \(C - \{0\}\) 是无尖锥。

若 C 是无尖凸锥，则 \(C \cup \{0\}\) 是尖凸锥。然而，若 C 是尖凸锥，\(C - \{0\}\) 不一定是凸的。我们称尖凸锥为真尖锥，如果它不包含任何过 0 的直线。于是

命题 9. --- 尖凸锥 C 是真尖锥，当且仅当在 C 中对 0 取补所得的无尖锥 C' 是凸的。

若 C 包含一条过 0 的直线，则显然 \(C'\) 不是凸的。现在设 C 是真尖锥，并设 x、y 是 \(C'\) 的两点。以 x、y 为端点的闭线段包含在 C 中；若它包含 0，则 \(\lambda x + (1 - \lambda)y = 0\)（对某个 \(\lambda\)，满足 \(0 < \lambda < 1\)），因此 \(x = \mu y\)，其中 \(\mu < 0\)。于是 C 包含过 0 与 x 的直线，与假设矛盾。

命题 10. --- E 的子集 C 是凸锥，当且仅当 \(C + C \subset C\) 且 \(\lambda C \subset C\)（对所有 \(\lambda > 0\)）。

因为条件「\(\lambda C \subset C\)（对所有 \(\lambda > 0\)）」刻画了锥。若 C 是凸的，我们有 \(C + C = \frac{1}{2}C + \frac{1}{2}C = C\)（II, 第 8 页, 注记）。反之，若锥 C 满足 \(C + C \subset C\)，则当 \(0 < \lambda < 1\) 时，\(\lambda C + (1 - \lambda)C = C + C \subset C\)，这表明 C 是凸的。

推论 1. --- 若 C 是非空凸锥，则 C 生成的向量空间是集 C - C（由 x - y（其中 x、y 在 C 中变化）组成的点集）。

因为，若 V = C - C，则 V 非空，我们有 \(\lambda V = V\)（对所有 \(\lambda \neq 0\)），且 \(V + V = C + C - (C + C) \subset C - C = V\)，这表明 V 是向量子空间。最后，每个包含 C 的向量子空间也包含 V。

推论 2. --- 若 C 是尖凸锥，则包含在 C 中的最大向量子空间是集 \(C \cap (-C)\)。

因为，若 \(W = C \cap (-C)\)，则 W 非空且 \(\lambda W = W\)（对所有 \(\lambda \neq 0\)），又有

\[
\mathrm{W} + \mathrm{W} \subset (\mathrm{C} + \mathrm{C}) \cap (- (\mathrm{C} + \mathrm{C})) \subset \mathrm{C} \cap (- \mathrm{C}) = \mathrm{W},
\]

这表明 W 是向量子空间。显然，包含在 C 中的每个向量子空间也都包含在 W 中。

显然，若 f 是 E 到向量空间 F 中的线性映射，则 E 中凸锥 C 的像 \(f(\mathbf{C})\) 是 F 中的凸锥。E 中（以 0 为顶点的）凸锥的每个交是凸锥。对 E 的每个子集 A，包含 A 的诸凸锥（这样的锥存在，E 本身就是一个）之交是包含 A 的最小凸锥；它称为由 A 生成的凸锥。

命题 11. --- 设 \((\mathbf{C}_{\iota})_{\iota \in \mathbf{I}}\) 是 E 中凸锥的一个族；由诸 \(\mathbf{C}_{\iota}\) 之并生成的凸锥与点 \(\sum_{\iota \in \mathbf{J}} x_{\iota}\) 组成的集恒同，其中 \(\mathbf{J}\) 是 I 的任意有限

子集，且 \(x_{\iota} \in \mathbf{C}_{\iota}\)（对所有 \(\iota \in \mathbf{J}\)）。

事实上，显然这样的点组成的集 C 是包含诸 \(C_{\iota}\) 之并的一个凸锥，且它包含在包含该并的任何凸锥中。

推论. --- 对 E 的任意子集 A，由 A 生成的凸锥与线性组合 \(\sum_{i\in J}\lambda_{i}x_{i}\) 组成的集恒同，其中 \((x_{i})_{i\in J}\) 是 A 中点的任意有限非空族，且 \(\lambda_{i}>0\)（对所有 \(i\in J\)）。

只需看到：若一个凸锥包含 A 的一点 \(x \neq 0\)，则它也包含半直线 \(C_{x}\)（由点 \(\lambda x\)（\(\lambda\) 在正数集中变化）组成），并且 \(C_{x}\) 是凸锥。

命题 12. --- 若 A 是 E 中的凸集，则由 A 生成的凸锥与 \(C = \bigcup_{\lambda > 0} \lambda A\) 恒同。

集 C 显然是锥；只需证明 C 是凸的。设 \(\lambda x, \mu y\) 是 C 的两点（\(\lambda > 0, \mu > 0, x \in A, y \in A\)）。设 \(\alpha, \beta\) 是两个 \textgreater{} 0 的数，满足 \(\alpha + \beta = 1\)。则 \(\alpha \lambda x + \beta \mu y = (\alpha \lambda + \beta \mu) z\)，其中 \(z \in A\)，且 \(\alpha \lambda + \beta \mu > 0\)；故 \(\alpha \lambda x + \beta \mu y \in C\)。

注记. --- 1) 在 prop. 12 的假设下，若 \(0 \notin \mathbf{A}\)，则锥 \(C\) 是无尖锥，从而 \(C \cup \{0\}\) 是真尖锥。

\begin{enumerate}
\def\labelenumi{\arabic{enumi})}
\setcounter{enumi}{1}
\tightlist
\item
  设 A 是 E 中任意凸集；考察凸集 \(A_{1} = A \times \{1\}\)（在空间 \(F = E \times R\) 中）以及由 \(A_{1}\) 生成的以 0 为顶点的凸锥 C。Prop. 12 表明，\(A_{1}\) 是 C 与 F 中超平面 \(E \times \{1\}\) 的交。因此，E 中每个凸集都可以看作 F 中一个以 0 为顶点的凸锥与超平面 \(E \times \{1\}\) 的交在 E 上的投影。
\end{enumerate}

